\documentclass[10pt, letterpaper]{article}

\usepackage[
    top=2cm,
    bottom=2cm,
    left=2cm,
    right=2cm
]{geometry}
\usepackage{titlesec}
\usepackage{tabularx}
\usepackage{enumitem}
\usepackage[dvipsnames]{xcolor}
\definecolor{primaryColor}{RGB}{0, 79, 144}
\usepackage{fontawesome5}
\usepackage{paracol}
\usepackage{hyperref}
\hypersetup{
    colorlinks=true,
    urlcolor=primaryColor
}
\usepackage[french]{babel}
\usepackage{datetime2}
\usepackage{lastpage}
\usepackage{fancyhdr}
\usepackage{lmodern}

% ----- PIED DE PAGE -----
\pagestyle{fancy}
\fancyhf{}
\fancyfoot[R]{\small\textit{\color{gray}Dernière mise à jour le \DTMtoday}}
\fancyfoot[L]{\small\textit{\color{gray}Charly Romy TANGA – Page \thepage\ sur \pageref*{LastPage}}}
\renewcommand{\headrulewidth}{0pt}
\renewcommand{\footrulewidth}{0pt}

% ----- FORMAT DES SECTIONS -----
\titleformat{\section}{\large\bfseries\color{primaryColor}}{}{0pt}{}
\titlespacing*{\section}{0pt}{0.2cm}{0.15cm}

\setlist[itemize]{leftmargin=0.5cm, itemsep=3pt}

% ----- DOCUMENT -----
\begin{document}

% ===== EN-TÊTE =====
\begin{center}
    {\Large \textbf{Charly-Romy TANGA}}\\[3pt]
    {\small Paris, France \quad$\cdot$\quad 
    \href{mailto:charlyromytanga.cytech@gmail.com}{charlyromytanga.cytech@gmail.com} \quad$\cdot$\quad
    Tél : +33 6 15 42 25 74}\\[3pt]
    \href{https://linkedin.com/in/charly-romy.tanga}{\faLinkedinIn\ charly-romy.tanga} \quad$\cdot$\quad
    \href{https://github.com/charlyromytanga}{\faGithub\ charlyromytanga}
\end{center}

\vspace{0.2cm}

% ===== INTRO =====
\section*{Porté par la collaboration et les défis techniques}
Passionné par l’application des mathématiques et de l’informatique à la finance, je vise à développer et optimiser des modèles quantitatifs et des outils d’analyse soutenant les activités de trading et de gestion des risques. Mon objectif est de contribuer à des équipes pluridisciplinaires en concevant du code efficace, évolutif et bien documenté en \textbf{Python} et \textbf{Java} pour des solutions financières fondées sur les données.

% ===== EXPÉRIENCES =====
\section*{Expériences}

\textbf{Alternant Data Analyst – NaTran France (Spécialiste national des données gazières)} \hfill \textit{Septembre 2023 – Présent}\\
\textit{Ingénierie de la données, Département Data}
\begin{itemize}
    \item Développement de tableaux de bord nationaux sous \textbf{Power BI}, intégrant des jeux de données hétérogènes et réduisant le temps de reporting de 40\%.
    \item Conception et optimisation de pipelines de traitement de données avec \textbf{Python (Pandas, PySpark)}, automatisant des tâches récurrentes et réduisant le travail manuel de 70\%.
    \item Création et maintenance de modèles de données SQL pour les rapports analytiques, améliorant la qualité et l’accessibilité des données.
    \item Contribution à un projet web de visualisation de données (\textbf{Angular/JavaScript}), améliorant la réactivité du système de 30\%.
\end{itemize}

\vspace{0.2cm}

\textbf{Projet académique – Simulation en finance quantitative (CY Tech)} \hfill \textit{2025 (en cours)}\\
\textit{Projet personnel, Modélisation quantitative}
\begin{itemize}
    \item Implémentation d’un \textbf{pricer C++/Python} pour les options européennes à l’aide du modèle \textbf{Black-Scholes}.
    \item Simulation de processus stochastiques et utilisation de méthodes de \textbf{Monte Carlo} pour estimer des indicateurs de risque tels que la \textbf{Value-at-Risk (VaR)}.
    \item Application de méthodes numériques et d’optimisations de performance pour un calcul évolutif.
\end{itemize}

% ===== COMPÉTENCES ET TECHNOLOGIES =====
\section*{Compétences et Technologies}

\textbf{Langages de programmation :} Python, Java, C, SQL\\
\textbf{Bibliothèques et outils :} NumPy, pandas, SciPy, scikit-learn, PySpark, Git, Docker\\
\textbf{Mathématiques et finance :} Probabilités, Statistiques, Processus stochastiques, Algèbre linéaire, Modélisation financière\\
\textbf{Gestion des données :} Bases SQL/NoSQL, ETL, Pipelines de données\\
\textbf{Pratiques de développement :} Programmation orientée objet, Tests, Débogage, Documentation\\
\textbf{Langues :} Français (natif), Anglais (C1), Allemand (B1)

% ===== CERTIFICATIONS PROFESSIONNELLES =====
\section*{Certifications Professionnelles}

\textbf{Microsoft Certified: Data Analyst Associate} (en cours) – Spécialisation Power BI\\
\textbf{Bloomberg Market Concepts (BMC)} (en cours) – Fondamentaux des marchés financiers\\
\textbf{Certification AMF} (en cours) – Réglementation et conformité financières

% ===== FORMATION =====
\section*{Formation}
\textbf{Master 2 en Mathématiques appliquées et Informatique – Spécialisation FinTech} \hfill \textit{CY Tech, Cergy – 2024–2026}\\
\textbf{Licence en Mathématiques et Informatique} \hfill \textit{CY Tech, Cergy – 2023–2024}\\
Cours pertinents : Probabilités, Statistiques, Processus stochastiques, Fouille de données, Finance computationnelle

\end{document}
